\chapter*{ABSTRACT}
\addcontentsline{toc}{chapter}{ABSTRACT}

Schedule and cost management in large-scale logistics projects face high uncertainty and require solving the Resource-Constrained Project Scheduling Problem (RCPSP). Traditional network analysis methods (CPM/PERT) often overlook the dynamic nature of field operations, while exact mathematical solvers (e.g., CP-SAT) suffer from combinatorial explosion (NP-hard complexity) when applied to complex project structures.

To address this gap, this thesis proposes \textbf{GLPO (Graph-Based Logistics Project Optimizer)} - an end-to-end system combining Artificial Intelligence and Combinatorial Optimization. Unlike previous models, GLPO represents projects as directed heterogeneous graphs comprising 3 node types (Task, Resource, Shift) and 5 edge relations. The core of the system is a 3-stage AI Pipeline: (1) A Heterogeneous Graph Transformer (HGT) enhanced with Masked Autoencoder (MAE) self-supervised learning predicts delay risk distributions; (2) Monte Carlo CPM simulations (10,000 iterations) quantify uncertainties into Criticality Indices (CI); and (3) The Google OR-Tools (CP-SAT) solver utilizes CI as a learned heuristic to prune the search space, thereby rapidly converging to Pareto-optimal frontiers balancing Time, Cost, and Risk based on a customized 6-group financial model.

Empirical evaluations on the real-world C2012-04 dataset (67 tasks) demonstrate that the HGT model achieves high predictive accuracy (MAE = 0.91 days for expected delay). The integration of AI-driven heuristics enables the CP-SAT solver to discover optimal crashing scenarios and reduces the solution search time by 25.6\% compared to baseline mathematical programming.

Architecturally, the system is deployed as independent microservices with a Backend (FastAPI, PostgreSQL) and a Frontend (React, Vite) for network visualization, containerized via Docker. This establishes a comprehensive, data-driven, and highly adaptive project management tool.